\startcomponent ctx-estructuras
\product ppchtextut


\startchapter[title={Estructuras}]
  
  Vamos a familiarizarnos con algunas estructuras químicas.

  % ##############################################################################
  \startsection[title={Hexagonal}]
    Para empezar, utilizaremos aquí la estructura \type{SIX}, que utiliza como base
    un grupo de seis átomos formando un hexágono regular. Cada átomo está asociado
    a un número.
  
    En la siguiente imagen se aprecia la numeración de cada átomo del hexágono
    asociado a la estructura \type{SIX}
  
    % --------------------------------------------------------------------------
    % SIX
    % ---------------------------------------------------------------------------    
    % Definición de buffer
    % Para imagen
    \startbuffer[molecula]
      \startchemical[frame=on,width=4000,top=1200,bottom=1800]
        \bottext{SIX}
        \chemical[SIX,SB,Z][1,2,3,4,5,6]
      \stopchemical
    \stopbuffer
    % Para código fuente
    \startbuffer[moleculatxt]
      \startchemical
        \bottext{SIX}
        \setupchemical[width=4000,top=1200,bottom=1500]
        \chemical[SIX,SB,Z][1,2,3,4,5,6]
      \stopchemical
    \stopbuffer
    % .........................................................................
    % Definimos una estructura que alinee los centros
    \startlinealignment[left]
      \framed[frame=off,location=middle]{
        \getbuffer[molecula]}
      \hskip 0.5em % Espacio entre fórmula y código
      \framed[frame=off,location=middle,align=right]
             {\switchtobodyfont[8.0pt]\typebuffer[moleculatxt]}
    \stoplinealignment
    % ---------------------------------------------------------------------------
    A continuación, explicamos por encima el significado del código fuente utilizado
    para representar la figura anterior. Vemos que al comando \type{chemical}
    se le añaden dos grupos de argumentos.
    \startitemize[packed]
    \item El primer grupo contiene \type{SB} (\emph{single bonds}):
      \startitemize[packed]
      \item \type{SIX} organiza una estructura hexagonal.
      \item \type{SB} dibuja los enlaces (\type{B}) sencillos (\type{S})
        \emph{dejando espacio para dibujar los átomos} que se encuentran en los
        vértices de la estructura, que en este caso hemos aprovecado para poner
        el número que le corresponde ºa cada vértice. Si se hubiera puesto solo
        \type{B}, los enlaces unirían los vértices, no dejando que se vieran los
        átomos situados en los vértices.
      \stopitemize
    \item El segundo grupo contiene los átomos que se sitúan en la estructura
      hexagonal que, como hemos dicho son los números que les corresponden,
      aunque podrían ser \chemical{C} o cualquier otro átomo que interese.
    \item El comando \type{\bottext} escribe en este caso \type{SIX} debajo de
      la fórmula (para ello, hemos aumentado el comando bottom en
      \type{\startchemical} para dejar el espacio apropiado).
    \stopitemize
    % ****************************************************************************
    \startsubsection[title={Fórmula esquemática del ciclohexano}]
      
      En el primer ejemplo veremos la fórmula esquemática del ciclohexano, donde las
      líneas representan los seis enlaces \type{C-C}, los vértices representan los
      átomos de carbono y los hidrógenos no aparecen.
      
      Vemos que al comando \type{\chemical} se le añaden inmediatamente un grupo de
      argumentos, que son \type{[SIX,B]}. La opción \type{SIX} utiliza la estructura
      formada por seis átomos formando un hexágono regular. La otra opción \type{B},
      dibuja todos los enlaces entre átomos contiguos, de vértice a vértice,
      \emph{sin dejar espacio a los átomos de la estructura}\/.
      
      En cuanto al comando \type{\startchemical},  opción \type{top=1200} sube (o
      baja) el límite superior del marco \type{frame=on} un cierto espacio en
      blanco; por otro lado, \type{bottom=1200} baja el límite inferior del marco,
      dejando un cierto espacio en blanco. En el resto de imágnes obviaremos el
      argumento \type{frame=on} en la sección derecha (texto), para que no se
      alargue el comando, aunque, cuando la fórmula aparezca recuadrada, el comando
      \type{frame=on} debe añadirse
    
      % --------------------------------------------------------------------------
      % Ciclohexano (fórmula esquemática)
      % --------------------------------------------------------------------------
      % Definición de buffer
      % Para imagen
      \startbuffer[molecula]
        \startchemical[frame=on,width=4000,top=1200,bottom=1500]
          \bottext{ciclohexano}
          \chemical[SIX,B]
        \stopchemical
      \stopbuffer
      % Para código fuente
      \startbuffer[moleculatxt]
        \startchemical
          \bottext{ciclohexano}
          \setupchemical[width=4000,top=1200,bottom=1500]
          \chemical[SIX,B]
        \stopchemical
      \stopbuffer
      % .........................................................................
      % Definimos una estructura que alinee los centros
      \startlinealignment[left]
        \framed[frame=off,location=middle]{
          \getbuffer[molecula]}
        \hskip 0.5em % Espacio entre fórmula y código
        \framed[frame=off,location=middle,align=right]
               {\switchtobodyfont[8.0pt]\typebuffer[moleculatxt]}
      \stoplinealignment
      % ---------------------------------------------------------------------------
    \stopsubsection
    
    % ***************************************************************************
    \startsubsection[title={Fórmula esquemática del butano}]
    
      Se podrían dibujar menos enlaces en esa estructura básica hexagonal.
      Por ejemplo, el butano
  
      % --------------------------------------------------------------------------
      % Butano (fórmula esquemática 1)
      % --------------------------------------------------------------------------
      % Definición de buffer
      % Para imagen
      \startbuffer[molecula]
        \startchemical[frame=on,left=1800,width=4000,top=1000,bottom=1500]
          \bottext{butano}
          \chemical[SIX,B612]
        \stopchemical
      \stopbuffer
      % Para código fuente
      \startbuffer[moleculatxt]
        \startchemical
          \bottext{butano}
          \setupchemical[left=1800,width=4000,top=1000,bottom=1500]
          \chemical[SIX,B612]
        \stopchemical
      \stopbuffer
      % .........................................................................
      % Definimos una estructura que alinee los centros
      \startlinealignment[left]
        \framed[frame=off,location=middle]{
          \getbuffer[molecula]}
        \hskip 0.5em % Espacio entre fórmula y código
        \framed[frame=off,location=middle,align=right]
               {\switchtobodyfont[8.0pt]\typebuffer[moleculatxt]}
      \stoplinealignment
      % ---------------------------------------------------------------------------
  
      Cada número que se le pasa a la opción \type{B} (enlaces) indica que se
      debe dibujar un enlace (de vértice a vértice) empezando por el número que
      se indica, hasta el siguiente átomo en la numeración. Así, en esta fórmula,
      \type{B612} indica que se dibujen los enlaces, del 6 al 1,  del 1 al 2 y, por
      último, del 2 al 3. Nótese que se podrían haber especificado los números en
      otro orden, por ejemplo \type{B261}.

      % --------------------------------------------------------------------------
      % Butano (fórmula esquemática 2)
      % --------------------------------------------------------------------------
      % Definición de buffer
      % Para imagen
      \startbuffer[molecula]
        \startchemical[frame=on,left=1800,width=4000,top=1300,bottom=600]
          \bottext{butano}
          \chemical[SIX,B56,R1]
        \stopchemical
      \stopbuffer
      % Para código fuente
      \startbuffer[moleculatxt]
        \startchemical
          \bottext{butano}
          \setupchemical[left=1800,width=4000,top=1000,bottom=1500]
          \chemical[SIX,B612]
        \stopchemical
      \stopbuffer
      % .........................................................................
      % Definimos una estructura que alinee los centros
      \startlinealignment[left]
        \framed[frame=off,location=middle]{
          \getbuffer[molecula]}
        \hskip 0.5em % Espacio entre fórmula y código
        \framed[frame=off,location=middle,align=right]
               {\switchtobodyfont[8.0pt]\typebuffer[moleculatxt]}
      \stoplinealignment
      % ---------------------------------------------------------------------------
      
    \stopsubsection
  
    % ****************************************************************************
    \startsubsection[title={Fórmula esquemática del metilciclohexano}]

      En esta fórmula añadimos la opción \type{+SR1}, que dibuja un radical acortado
      al final, es decir, sale del vértice 1 y llega al final acortado, bien porque
      se quiera así, o porque al final del mismo queramos escribir un átomo que,
      en caso contrario sería tachado por el radical. También podríamos haber
      obtenido el mismo efecto poniendo \type{+R1}, pero en el manual prefieren
      ponerlo como \type{+SR1}.
    
      % --------------------------------------------------------------------------
      % Metilciclohexano (fórmula esquemática)
      % ---------------------------------------------------------------------------
      % Definición de buffer
      % Para imagen
      \startbuffer[molecula]
        \startchemical[frame=on,width=4200,left=1900,top=1200,bottom=1500]
          \bottext{metilciclohexano}
          \chemical[SIX,B,+SR1]
        \stopchemical
      \stopbuffer
      % Para código fuente
      \startbuffer[moleculatxt]
        \startchemical
          \bottext{metilciclohexano}
          \setupchemical[width=4200,left=1900,top=1200,bottom=1500]
          \chemical[SIX,B,+SR1]
        \stopchemical
      \stopbuffer
      % .........................................................................
      % Definimos una estructura que alinee los centros
      \startlinealignment[left]
        \framed[frame=off,location=middle]{
          \getbuffer[molecula]}
        \hskip 0.5em % Espacio entre fórmula y código
        \framed[frame=off,location=middle,align=right]
               {\switchtobodyfont[8.0pt]\typebuffer[moleculatxt]}
      \stoplinealignment
      % ---------------------------------------------------------------------------
    
      Conviene tener en cuenta las siguientes opciones:
      \startitemize[packed,joinedup]
      \item \[\type{R}]: Radical sencillo largo (sale del vértice).
        No se acorta ni al principio ni al final. Se utiliza cuando nos interesa
        esa longitud. No se debe añadir ningún átomo al final porque quedaría
        tapado.
      \item \[\type{+SR}\] o [\type{+R}]: Radical acortado al final (sale del
        vértice).
        Se utiliza cuando nos interesea un radical más corto que en el caso anterior,
        o cuando al final del mismo queremos poner un átomo o grupo de átomos, de
        manera que no se vean tachados por el radical.
      \item \[\type{-SR}\] o [\type{-R}]: Radical acortado al principio
        (no sale del vértice), pero tiene el alcance del radical \type{R} al final.
        Se utiliza cuando se desea un radical largo, que no requiera un átomo al
        final, pero que en el vértice aparezca un átomo.
      \item \[\type{SR}\]: Radical acortado al principio y al final,
        (no sale del vértice), pero es más corto que el anterior. Se utiliza cuando
        queremos poner un átomo en el vértice de la estructura y nos interesa un
        radical corto al final, o bien queremos poner otro átomo al final.
      \stopitemize
    \stopsubsection
    
    % ****************************************************************************
    \startsubsection[title={Fórmula mixta del metilciclohexano}]
    
      % --------------------------------------------------------------------------
      % Metilciclohexano (fórmula mixta)
      % --------------------------------------------------------------------------
      % Definición de buffer
      % Para imagen
      \startbuffer[molecula]
        \startchemical[frame=on,width=4200,left=1900,top=1200,bottom=1500]
          \bottext{metilciclohexano}
          \chemical[SIX,B,+SR1,RZ1][CH_3]
        \stopchemical
      \stopbuffer
      % Para código fuente
      \startbuffer[moleculatxt]
        \startchemical
          \bottext{metilciclohexano}
          \setupchemical[width=4000,top=1200,bottom=1500]
          \chemical[SIX,B,+SR1,RZ1]{CH_3}
        \stopchemical
      \stopbuffer
      % .........................................................................
      % Definimos una estructura que alinee los centros
      \startlinealignment[left]
        \framed[frame=off,location=middle]{
          \getbuffer[molecula]}
        \hskip 0.5em % Espacio entre fórmula y código
        \framed[frame=off,location=middle,align=right]
               {\switchtobodyfont[8.0pt]\typebuffer[moleculatxt]}
      \stoplinealignment
      % ****************************************************************************
    
      % --------------------------------------------------------------------------
      % 1,2-dimetil-ciclohexano
      % --------------------------------------------------------------------------
      % Definición de buffer
      % Para imagen
      \startbuffer[molecula]
        \startchemical[frame=on,width=4000,left=1500,top=1200,bottom=1700]
          \bottext{1,2-dimetilclohexano}
          \chemical[SIX,B,+SR12,RZ12][CH_3,CH_3]
        \stopchemical
      \stopbuffer
      % Para código fuente
      \startbuffer[moleculatxt]
        \startchemical
          \bottext{1,2-dimetilciclohexano}
          \setupchemical[width=4000,left=1500,top=1200,bottom=1700]
          \chemical[SIX,B,+SR12,RZ12][CH_3,CH_3]
        \stopchemical
      \stopbuffer
      % .........................................................................
      % Definimos una estructura que alinee los centros
      \startlinealignment[left]
        \framed[frame=off,location=middle]{
          \getbuffer[molecula]}
        \hskip 0.5em % Espacio entre fórmula y código
        \framed[frame=off,location=middle,align=right]
               {\switchtobodyfont[8.0pt]\typebuffer[moleculatxt]}
      \stoplinealignment

    \stopsubsection
    
%      % --------------------------------------------------------------------------
%      % Ciclohexano bote
%      % --------------------------------------------------------------------------
%      % Definición de buffer
%      % Para imagen
%      \startbuffer[molecula]
%        \startchemical[frame=on,width=4000,top=1200,bottom=1600]
%          \bottext{Ciclohexano (bote)}
%          \chemical[BOAT,BB]
%        \stopchemical
%      \stopbuffer
%      % Para código fuente
%      \startbuffer[moleculatxt]
%        \startchemical
%          \uptext{(bote)}
%          \bottext{ciclohexano (bote)}
%          \setupchemical[width=4000,top=1200,bottom=1600]
%          \chemical[BOAT,BB]
%        \stopchemical
%      \stopbuffer
%      % .........................................................................
%      % Definimos una estructura que alinee los centros
%      \startlinealignment[left]
%        \framed[frame=off,location=middle]{
%          \getbuffer[molecula]}
%      \hskip 0.5em % Espacio entre fórmula y código
%      \framed[frame=off,location=middle,align=right]
%             {\switchtobodyfont[8.0pt]\typebuffer[moleculatxt]}
%      \stoplinealignment
%
%      % --------------------------------------------------------------------------
%      % Ciclohexano silla
%      % --------------------------------------------------------------------------
%      % Definición de buffer
%      % Para imagen
%      \startbuffer[molecula]
%        \startchemical[frame=on,width=4000,top=1200,bottom=1500]
%          \bottext{ciclohexano (silla)}
%          \chemical[CHAIR,BB]
%        \stopchemical
%      \stopbuffer
%      % Para código fuente
%      \startbuffer[moleculatxt]
%        \startchemical
%          \uptext{(bote)}
%          \bottext{ciclohexano (silla)}
%          \setupchemical[width=4000,top=1200,bottom=1500]
%          \chemical[CHAIR,BB]
%        \stopchemical
%      \stopbuffer
%      % .........................................................................
%      % Definimos una estructura que alinee los centros
%      \startlinealignment[left]
%        \framed[frame=off,location=middle]{
%          \getbuffer[molecula]}
%        \hskip 0.5em % Espacio entre fórmula y código
%        \framed[frame=off,location=middle,align=right]
%               {\switchtobodyfont[8.0pt]\typebuffer[moleculatxt]}
%      \stoplinealignment
%
%    \stopsubsection

  \stopsection

  % ##############################################################################
  \startsection[title={Estructura con átomo central}]
    Para empezar, utilizaremos aquí la estructura \type{ONE}, que utiliza como
    base un grupo de seis átomos formando un hexágono regular. Cada átomo está
    asociado a un número.
  
    Numeración de cada átomo asociado a la estructura \type{ONE}
 
    % --------------------------------------------------------------------------
    % ONE
    % ---------------------------------------------------------------------------    
    % Definición de buffer
    % Para imagen
    \startbuffer[molecula]
      \startchemical[frame=on,width=4000,top=1200,bottom=1800]
        \bottext{ONE}
        \chemical[ONE,SB,Z0,Z][A,1,2,3,4,5,6,7,8]
      \stopchemical
    \stopbuffer
    % Para código fuente
    \startbuffer[moleculatxt]
      \startchemical
        \bottext{ONE}
        \setupchemical[width=4000,top=1200,bottom=1500]
        \chemical[ONE,SB,Z0,Z][A,1,2,3,4,5,6,7,8]
      \stopchemical
    \stopbuffer
    % .........................................................................
    % Definimos una estructura que alinee los centros
    \startlinealignment[left]
      \framed[frame=off,location=middle]{
        \getbuffer[molecula]}
      \hskip 0.5em % Espacio entre fórmula y código
      \framed[frame=off,location=middle,align=right]
             {\switchtobodyfont[8.0pt]\typebuffer[moleculatxt]}
    \stoplinealignment
    % ---------------------------------------------------------------------------

    Vemos que el átomo central se numera con cero \type{Z0}. Los átomos unidos
    a este se numeran del uno al ocho \type{Z1...Z8}.
  \stopsubsection

  % ****************************************************************************
  \startsubsection[title={Fórmula desarrollada del metano}]
    
    % --------------------------------------------------------------------------
    % Metilciclohexano (fórmula mixta)
    % --------------------------------------------------------------------------
    % Definición de buffer
    % Para imagen
    \startbuffer[molecula]
      \startchemical[frame=on,width=4200,left=1900,top=1300,bottom=1800]
        \bottext{metano}
        \chemical[ONE,SR1357,Z0,Z1357][C,H,H,H,H]
      \stopchemical
    \stopbuffer
    % Para código fuente
    \startbuffer[moleculatxt]
      \startchemical
        \bottext{metano}
        \setupchemical[width=4000,top=1300,bottom=1800]
        \chemical[ONE,SR1357,Z0,Z1357]{C,H,H,H,H}
      \stopchemical
    \stopbuffer
    % .........................................................................
    % Definimos una estructura que alinee los centros
    \startlinealignment[left]
      \framed[frame=off,location=middle]{
        \getbuffer[molecula]}
      \hskip 0.5em % Espacio entre fórmula y código
      \framed[frame=off,location=middle,align=right]
             {\switchtobodyfont[8.0pt]\typebuffer[moleculatxt]}
           \stoplinealignment
  % ***************************************************************************
           
      Nótese \type{SR} que acorta los enlaces al principio y al final.   
    \stopsubsection
  \stopsection
  
  % ##############################################################################
  \startsection[title=Estructura tiangular]
    Para empezar, utilizaremos aquí la estructura \type{THREE}, que utiliza como
    base un grupo de seis átomos formando un hexágono regular. Cada átomo está
    asociado a un número.
  
    Numeración de cada átomo del pentágono asociado a la estructura \type{THREE}
 
    % --------------------------------------------------------------------------
    % THREE
    % ---------------------------------------------------------------------------    
    % Definición de buffer
    % Para imagen
    \startbuffer[molecula]
      \startchemical[frame=on,width=4000,top=1200,bottom=1800]
        \bottext{THREE}
        \chemical[THREE,SB,Z][1,2,3]
      \stopchemical
    \stopbuffer
    % Para código fuente
    \startbuffer[moleculatxt]
      \startchemical
        \bottext{THREE}
        \setupchemical[width=4000,top=1200,bottom=1500]
        \chemical[THREE,SB,Z][1,2,3]
      \stopchemical
    \stopbuffer
    % .........................................................................
    % Definimos una estructura que alinee los centros
    \startlinealignment[left]
      \framed[frame=off,location=middle]{
        \getbuffer[molecula]}
      \hskip 0.5em % Espacio entre fórmula y código
      \framed[frame=off,location=middle,align=right]
             {\switchtobodyfont[8.0pt]\typebuffer[moleculatxt]}
    \stoplinealignment
    % ---------------------------------------------------------------------------

  \stopsection

  % ##############################################################################
  \startsection[title=Estructura cuadrangular]
    Para empezar, utilizaremos aquí la estructura \type{FOUR}, que utiliza como
    base un grupo de seis átomos formando un hexágono regular. Cada átomo está
    asociado a un número.
  
    Numeración de cada átomo del pentágono asociado a la estructura \type{FOUR}
 
    % --------------------------------------------------------------------------
    % FOUR
    % ---------------------------------------------------------------------------    
    % Definición de buffer
    % Para imagen
    \startbuffer[molecula]
      \startchemical[frame=on,width=4000,top=1200,bottom=1500]
        \bottext{FOUR}
        \chemical[FOUR,SB,Z][1,2,3,4]
      \stopchemical
    \stopbuffer
    % Para código fuente
    \startbuffer[moleculatxt]
      \startchemical
        \bottext{FOUR}
        \setupchemical[width=4000,top=1200,bottom=1500]
        \chemical[FOUR,SB,Z][1,2,3,4]
      \stopchemical
    \stopbuffer
    % .........................................................................
    % Definimos una estructura que alinee los centros
    \startlinealignment[left]
      \framed[frame=off,location=middle]{
        \getbuffer[molecula]}
      \hskip 0.5em % Espacio entre fórmula y código
      \framed[frame=off,location=middle,align=right]
             {\switchtobodyfont[8.0pt]\typebuffer[moleculatxt]}
    \stoplinealignment
    % ---------------------------------------------------------------------------

  \stopsection
  
  % ##############################################################################
  \startsection[title=Estructura pentagonal]
    Para empezar, utilizaremos aquí la estructura \type{FIVE}, que utiliza como
    base un grupo de seis átomos formando un hexágono regular. Cada átomo está
    asociado a un número.
  
    Numeración de cada átomo del pentágono asociado a la estructura \type{FIVE}
 
    % --------------------------------------------------------------------------
    % FIVE
    % ---------------------------------------------------------------------------    
    % Definición de buffer
    % Para imagen
    \startbuffer[molecula]
      \startchemical[frame=on,width=4000,top=1200,bottom=1700]
        \bottext{FIVE}
        \chemical[FIVE,SB,Z][1,2,3,4,5]
      \stopchemical
    \stopbuffer
    % Para código fuente
    \startbuffer[moleculatxt]
      \startchemical
        \bottext{FIVE}
        \setupchemical[width=4000,top=1100,bottom=1700]
        \chemical[FIVE,SB,Z][1,2,3,4,5]
      \stopchemical
    \stopbuffer
    % .........................................................................
    % Definimos una estructura que alinee los centros
    \startlinealignment[left]
      \framed[frame=off,location=middle]{
        \getbuffer[molecula]}
      \hskip 0.5em % Espacio entre fórmula y código
      \framed[frame=off,location=middle,align=right]
             {\switchtobodyfont[8.0pt]\typebuffer[moleculatxt]}
    \stoplinealignment
    % ---------------------------------------------------------------------------

    \startsubsection[title=ciclopentano]
    % --------------------------------------------------------------------------
    % Ciclopentano (fórmula esquemática)
    % --------------------------------------------------------------------------
    % Definición de buffer
    % Para imagenn
    \startbuffer[molecula]
      \startchemical[frame=on,width=4000,top=1200,bottom=1300]
        \bottext{ciclopentano}
        \chemical[FIVE,B]
      \stopchemical
    \stopbuffer
    % Para código fuente
    \startbuffer[moleculatxt]
      \startchemical
        \bottext{ciclopentano}
        \setupchemical[width=4000,top=1200,bottom=1500]
        \chemical[FIVE,B]
      \stopchemical
    \stopbuffer
    % .........................................................................
    % Definimos una estructura que alinee los centros
    \startlinealignment[left]
      \framed[frame=off,location=middle]{
        \getbuffer[molecula]}
      \hskip 0.5em % Espacio entre fórmula y código
      \framed[frame=off,location=middle,align=right]
             {\switchtobodyfont[8.0pt]\typebuffer[moleculatxt]}
    \stoplinealignment
    % ---------------------------------------------------------------------------
   \stopsubsection

   \startsubsection[title={\emph{Tetrahidrofurano}, oxolano u oxaciclopentano}]
    % --------------------------------------------------------------------------
    % Oxolano (fórmula esquemática)
    % --------------------------------------------------------------------------
    % Definición de buffer
    % Para imagenn
    \startbuffer[molecula]
      \startchemical[frame=on,width=4000,top=1200,bottom=1500]
        \bottext{\emph{tetrahidrofurano}}
        \chemical[FIVE,B451,+SB2,-SB3,Z3][O]
      \stopchemical
    \stopbuffer
    % Para código fuente
    \startbuffer[moleculatxt]
      \startchemical
        \bottext{oxolano}
        \setupchemical[width=4000,top=1200,bottom=1500]
        \chemical[FIVE,B451,+SB2,-SB3,Z3][O]
      \stopchemical
    \stopbuffer
    % .........................................................................
    % Definimos una estructura que alinee los centros
    \startlinealignment[left]
      \framed[frame=off,location=middle]{
        \getbuffer[molecula]}
      \hskip 0.5em % Espacio entre fórmula y código
      \framed[frame=off,location=middle,align=right]
             {\switchtobodyfont[8.0pt]\typebuffer[moleculatxt]}
    \stoplinealignment
  % ---------------------------------------------------------------------------
  \stopsubsection

  \startsubsection[title=3-metiloxolano]
    % --------------------------------------------------------------------------
    % 3-metiloxolano (fórmula esquemática)
    % --------------------------------------------------------------------------
    % Definición de buffer
    % Para imagenn
    \startbuffer[molecula]
      \startchemical[frame=on,left=2000,width=4500,top=1200,bottom=1500]
        \bottext{3-metiloxolano}
        \chemical[FIVE,B145,+SB2,-SB3,Z3,+SR1,RZ1][O,CH_3]
      \stopchemical
    \stopbuffer
    % Para código fuente
    \startbuffer[moleculatxt]
      \startchemical
        \bottext{3-metiloxolano}
        \setupchemical[width=4000,top=1200,bottom=1500]
        \chemical[FIVE,B145,+SB2,-SB3,Z3,RZ1][O,CH_3]
      \stopchemical
    \stopbuffer
    % .........................................................................
    % Definimos una estructura que alinee los centros
    \startlinealignment[left]
      \framed[frame=off,location=middle]{
        \getbuffer[molecula]}
      \hskip 0.5em % Espacio entre fórmula y código
      \framed[frame=off,location=middle,align=right]
             {\switchtobodyfont[8.0pt]\typebuffer[moleculatxt]}
    \stoplinealignment
  % ---------------------------------------------------------------------------
  \stopsubsection
\stopsection

% ##############################################################################
\startsection[title=Estructura heptagonal]
  Para empezar, utilizaremos aquí la estructura \type{SEVEN}, que utiliza como
  base un grupo de seis átomos formando un hexágono regular. Cada átomo está
  asociado a un número.
  
  Numeración de cada átomo del pentágono asociado a la estructura \type{SEVEN}
 
  % --------------------------------------------------------------------------
  % SEVEN
  % ---------------------------------------------------------------------------    
  % Definición de buffer
  % Para imagen
  \startbuffer[molecula]
    \startchemical[frame=on,width=4000,top=1200,bottom=1800]
      \bottext{SEVEN}
      \chemical[SEVEN,SB,Z][1,2,3,4,5,6,7]
      \stopchemical
    \stopbuffer
    % Para código fuente
    \startbuffer[moleculatxt]
      \startchemical
        \bottext{SEVEN}
        \setupchemical[width=4000,top=1200,bottom=1500]
        \chemical[SEVEN,SB,Z][1,2,3,4,5,6,7]
      \stopchemical
    \stopbuffer
    % .........................................................................
    % Definimos una estructura que alinee los centros
    \startlinealignment[left]
      \framed[frame=off,location=middle]{
        \getbuffer[molecula]}
      \hskip 0.5em % Espacio entre fórmula y código
      \framed[frame=off,location=middle,align=right]
             {\switchtobodyfont[8.0pt]\typebuffer[moleculatxt]}
    \stoplinealignment
    % ---------------------------------------------------------------------------
  \stopsection

  % ##############################################################################
  \startsection[title=Estructura octogonal]
    Para empezar, utilizaremos aquí la estructura \type{EIGHT}, que utiliza como
    base un grupo de seis átomos formando un hexágono regular. Cada átomo está
    asociado a un número.
  
    Numeración de cada átomo del pentágono asociado a la estructura \type{EIGHT}
 
    % --------------------------------------------------------------------------
    % EIGHT
    % ---------------------------------------------------------------------------    
    % Definición de buffer
    % Para imagen
    \startbuffer[molecula]
      \startchemical[frame=on,width=4000,top=1200,bottom=2000]
        \bottext{EIGHT}
        \chemical[EIGHT,SB,Z][1,2,3,4,5,6,7,8]
      \stopchemical
    \stopbuffer
    % Para código fuente
    \startbuffer[moleculatxt]
      \startchemical
        \bottext{EIGHT}
        \setupchemical[width=4000,top=1200,bottom=2000]
        \chemical[EIGHT,SB,Z][1,2,3,4,5,6,7,8]
      \stopchemical
    \stopbuffer
    % .........................................................................
    % Definimos una estructura que alinee los centros
    \startlinealignment[left]
      \framed[frame=off,location=middle]{
        \getbuffer[molecula]}
      \hskip 0.5em % Espacio entre fórmula y código
      \framed[frame=off,location=middle,align=right]
             {\switchtobodyfont[8.0pt]\typebuffer[moleculatxt]}
    \stoplinealignment
    % ---------------------------------------------------------------------------

  \stopsection

  % ##############################################################################
  \startsection[title=Estructura nonagonal]
    Para empezar, utilizaremos aquí la estructura \type{NINE}, que utiliza como
    base un grupo de seis átomos formando un hexágono regular. Cada átomo está
    asociado a un número.
  
    Numeración de cada átomo del pentágono asociado a la estructura \type{NINE}
 
    % --------------------------------------------------------------------------
    % NINE
    % ---------------------------------------------------------------------------    
    % Definición de buffer
    % Para imagen
    \startbuffer[molecula]
      \startchemical[frame=on,width=4000,top=1200,bottom=2100]
        \bottext{NINE}
        \chemical[NINE,SB,Z][1,2,3,4,5,6,7,8]
      \stopchemical
    \stopbuffer
    % Para código fuente
    \startbuffer[moleculatxt]
      \startchemical
        \bottext{NINE}
        \setupchemical[width=4000,top=1200,bottom=2100]
        \chemical[NINE,SB,Z][1,2,3,4,5,6,7,8]
      \stopchemical
    \stopbuffer
    % .........................................................................
    % Definimos una estructura que alinee los centros
    \startlinealignment[left]
      \framed[frame=off,location=middle]{
        \getbuffer[molecula]}
      \hskip 0.5em % Espacio entre fórmula y código
      \framed[frame=off,location=middle,align=right]
             {\switchtobodyfont[8.0pt]\typebuffer[moleculatxt]}
    \stoplinealignment
    % ---------------------------------------------------------------------------

  \stopsection

  % ##############################################################################
  \startsection[title=Estructura nonagonal]
    Para empezar, utilizaremos aquí la estructura \type{NINE}, que utiliza como
    base un grupo de seis átomos formando un hexágono regular. Cada átomo está
    asociado a un número.
  
    Numeración de cada átomo del pentágono asociado a la estructura \type{NINE}
 
    % --------------------------------------------------------------------------
    % NINE
    % ---------------------------------------------------------------------------    
    % Definición de buffer
    % Para imagen
    \startbuffer[molecula]
      \startchemical[frame=on,width=4000,top=1200,bottom=2100]
        \bottext{NINE}
        \chemical[NINE,SB,Z][1,2,3,4,5,6,7,8]
      \stopchemical
    \stopbuffer
    % Para código fuente
    \startbuffer[moleculatxt]
      \startchemical
        \bottext{NINE}
        \setupchemical[width=4000,top=1200,bottom=2100]
        \chemical[NINE,SB,Z][1,2,3,4,5,6,7,8]
      \stopchemical
    \stopbuffer
    % .........................................................................
    % Definimos una estructura que alinee los centros
    \startlinealignment[left]
      \framed[frame=off,location=middle]{
        \getbuffer[molecula]}
      \hskip 0.5em % Espacio entre fórmula y código
      \framed[frame=off,location=middle,align=right]
             {\switchtobodyfont[8.0pt]\typebuffer[moleculatxt]}
    \stoplinealignment
    % ---------------------------------------------------------------------------

  \stopsection

  % ##############################################################################
  \startsection[title=Estructura hexagonal silla]
    Para empezar, utilizaremos aquí la estructura \type{CHAIR}, que utiliza como
    base un grupo de seis átomos. Cada átomo está asociado a un número.
    Hay dos tipos de enlace, numéricos y alfabéticos.
  
    Numeración de cada átomo ciclo asociado a la estructura \type{CHAIR}
 
    % --------------------------------------------------------------------------
    % CHAIR
    % ---------------------------------------------------------------------------    
    % Definición de buffer
    % Para imagen
    \startbuffer[molecula]
      \startchemical[frame=on,width=5000,top=1800,bottom=2200]
        \bottext{CHAIR}
        \chemical[CHAIR,B,+LSR,+RSR,+LRZ,+RRZ][a,b,c,d,e,f,1,2,3,4,5,6]
      \stopchemical
    \stopbuffer
    % Para código fuente
    \startbuffer[moleculatxt]
      \startchemical
        \bottext{CHAIR}
        \setupchemical[width=5000,top=1800,bottom=2200]
        \chemical
          [CHAIR,B,+LSR,+RSR,+LRZ,+RRZ]
          [a,b,c,d,e,f,1,2,3,4,5,6,7,8]
      \stopchemical
    \stopbuffer
    % .........................................................................
    % Definimos una estructura que alinee los centros
    \startlinealignment[left]
      \framed[frame=off,location=middle]{
        \getbuffer[molecula]}
      \hskip 0.5em % Espacio entre fórmula y código
      \framed[frame=off,location=middle,align=right]
             {\switchtobodyfont[8.0pt]\typebuffer[moleculatxt]}
    \stoplinealignment
    % ---------------------------------------------------------------------------
    
    \startsubsection[title={Ciclohexano (silla)}]
      % --------------------------------------------------------------------------
      % Ciclohexano silla
      % --------------------------------------------------------------------------
      % Definición de buffer
      % Para imagen
      \startbuffer[molecula]
        \startchemical[frame=on,width=4000,top=1200,bottom=1500]
          \bottext{ciclohexano (silla)}
          \chemical[CHAIR,B]
        \stopchemical
      \stopbuffer
      % Para código fuente
      \startbuffer[moleculatxt]
        \startchemical
          \uptext{(bote)}
          \bottext{ciclohexano (silla)}
          \setupchemical[width=4000,top=1200,bottom=1500]
          \chemical[CHAIR,B]
        \stopchemical
      \stopbuffer
      % .........................................................................
      % Definimos una estructura que alinee los centros
      \startlinealignment[left]
        \framed[frame=off,location=middle]{
          \getbuffer[molecula]}
        \hskip 0.5em % Espacio entre fórmula y código
        \framed[frame=off,location=middle,align=right]
               {\switchtobodyfont[8.0pt]\typebuffer[moleculatxt]}
      \stoplinealignment

      % --------------------------------------------------------------------------
      % Ciclohexano silla
      % --------------------------------------------------------------------------
      % Definición de buffer
      % Para imagen
      \startbuffer[molecula]
        \startchemical[frame=on,width=4000,top=1200,bottom=1500]
          \bottext{ciclohexano (silla)}
          \chemical[CHAIR,BB]
        \stopchemical
      \stopbuffer
      % Para código fuente
      \startbuffer[moleculatxt]
        \startchemical
          \uptext{(bote)}
          \bottext{ciclohexano (silla)}
          \setupchemical[width=4000,top=1200,bottom=1500]
          \chemical[CHAIR,BB]
        \stopchemical
      \stopbuffer
      % .........................................................................
      % Definimos una estructura que alinee los centros
      \startlinealignment[left]
        \framed[frame=off,location=middle]{
          \getbuffer[molecula]}
        \hskip 0.5em % Espacio entre fórmula y código
        \framed[frame=off,location=middle,align=right]
               {\switchtobodyfont[8.0pt]\typebuffer[moleculatxt]}
      \stoplinealignment
      
    \stopsubsection

  \stopsection

  % ##############################################################################
  \startsection[title=Estructura hexagonal bote]
    La estructura hexagonal bote utiliza como base un grupo de seis átomos. Cada
    átomo está asociado a un número.
    Hay dos tipos de enlace, numéricos y alfabéticos.
  
    Numeración de cada átomo ciclo asociado a la estructura \type{BOAT}
 
    % --------------------------------------------------------------------------
    % BOAT
    %---------------------------------------------------------------------------
    % Definición de buffer
    % Para imagen
    \startbuffer[molecula]
      \startchemical[frame=on,width=5000,top=1800,bottom=2200]
        \bottext{BOAT}
        \chemical[BOAT,B,+LSR,+RSR,+LRZ,+RRZ][a,b,c,d,e,f,1,2,3,4,5,6]
      \stopchemical
    \stopbuffer
    % Para código fuente
    \startbuffer[moleculatxt]
      \startchemical
        \bottext{BOAT}
        \setupchemical[width=5000,top=1800,bottom=2200]
        \chemical
          [BOAT,B,+LSR,+RSR,+LRZ,+RRZ]
          [a,b,c,d,e,f,1,2,3,4,5,6,7,8]
      \stopchemical
    \stopbuffer
    % .........................................................................
    % Definimos una estructura que alinee los centros
    \startlinealignment[left]
      \framed[frame=off,location=middle]{
        \getbuffer[molecula]}
      \hskip 0.5em % Espacio entre fórmula y código
      \framed[frame=off,location=middle,align=right]
             {\switchtobodyfont[8.0pt]\typebuffer[moleculatxt]}
    \stoplinealignment
    % ---------------------------------------------------------------------------
    
    \startsubsection[title={Ciclohexano (bote)}]
      % --------------------------------------------------------------------------
      % Ciclohexano bote
      % --------------------------------------------------------------------------
      % Definición de buffer
      % Para imagen
      \startbuffer[molecula]
        \startchemical[frame=on,width=4000,top=1200,bottom=1500]
          \bottext{ciclohexano (bote)}
          \chemical[BOAT,B]
        \stopchemical
      \stopbuffer
      % Para código fuente
      \startbuffer[moleculatxt]
        \startchemical
          \uptext{(bote)}
          \bottext{ciclohexano (bote)}
          \setupchemical[width=4000,top=1200,bottom=1500]
          \chemical[BOAT,B]
        \stopchemical
      \stopbuffer
      % .........................................................................
      % Definimos una estructura que alinee los centros
      \startlinealignment[left]
        \framed[frame=off,location=middle]{
          \getbuffer[molecula]}
        \hskip 0.5em % Espacio entre fórmula y código
        \framed[frame=off,location=middle,align=right]
               {\switchtobodyfont[8.0pt]\typebuffer[moleculatxt]}
      \stoplinealignment

      % --------------------------------------------------------------------------
      % Ciclohexano bote
      % --------------------------------------------------------------------------
      % Definición de buffer
      % Para imagen
      \startbuffer[molecula]
        \startchemical[frame=on,width=4000,top=1200,bottom=1500]
          \bottext{ciclohexano (bote)}
          \chemical[BOAT,BB]
        \stopchemical
      \stopbuffer
      % Para código fuente
      \startbuffer[moleculatxt]
        \startchemical
          \uptext{(bote)}
          \bottext{ciclohexano (bote)}
          \setupchemical[width=4000,top=1200,bottom=1500]
          \chemical[BOAT,BB]
        \stopchemical
      \stopbuffer
      % .........................................................................
      % Definimos una estructura que alinee los centros
      \startlinealignment[left]
        \framed[frame=off,location=middle]{
          \getbuffer[molecula]}
        \hskip 0.5em % Espacio entre fórmula y código
        \framed[frame=off,location=middle,align=right]
               {\switchtobodyfont[8.0pt]\typebuffer[moleculatxt]}
      \stoplinealignment
      
    \stopsubsection

  \stopsection

  
\stopchapter

\stopcomponent



  

%%% Local Variables:
%%% coding: utf-8
%%% mode: context
%%% ConTeXt-Mark-version: "IV"
%%% TeX-engine: default
%%% TeX-master: "../ctx-ppchtextut.tex"
%%% End:
